\section{Introduction}\label{sec:introduction}

The continuous Hochschild cohomology of an operator algebra measures
the obstruction to solving its bounded cocycle equations. The theory
also connects these equations with stability of multiplication under
small perturbations~\cite[Chapter~7]{SinclairSmith1995}.
For a von Neumann algebra acting on itself by multiplication, the
degree-one equation is the derivation equation, and its vanishing is
the inner-derivation theorem of Kadison and
Sakai~\cite{Kadison1966,Sakai1966}.
The higher-degree theory was developed by Johnson, Kadison, and
Ringrose~\cite{KadisonRingrose1971,KadisonRingrose1971II,JKR1972}.
Throughout, boundedness means ordinary
operator-norm boundedness of a multilinear map. No complete-boundedness
condition is imposed.

For a complex von Neumann algebra \(M\), let \(C_b^k(M,M)\) be the
Banach space of bounded complex \(k\)-linear maps \(M^k\to M\),
with norm
$\norm{f}=\sup_{\norm{a_1},\ldots,\norm{a_k}\leq1}
\norm{f(a_1,\ldots,a_k)}$, and put \(C_b^0(M,M)=M\).
For \(k\geq1\), the differential is
\begin{align}
(df)(a_1,\ldots,a_{k+1})
={}&a_1f(a_2,\ldots,a_{k+1})\notag\\
&+\sum_{j=1}^k(-1)^j
 f(a_1,\ldots,a_ja_{j+1},\ldots,a_{k+1})\notag\\
&+(-1)^{k+1}f(a_1,\ldots,a_k)a_{k+1}.
\label{eq:differential}
\end{align}
In degree zero, \((db)(a)=ab-ba\). As usual, \(d^2=0\).
We use the notation
\[
H_b^k(M,M)=
\ker\bigl(d:C_b^k(M,M)\to C_b^{k+1}(M,M)\bigr)
\big/
\operatorname{im}\bigl(d:C_b^{k-1}(M,M)\to C_b^k(M,M)\bigr).
\]
The denominator is the actual image, without a closure.

The Kadison--Ringrose cohomology conjecture asks whether
\(H_b^k(M,M)=0\) for every complex von Neumann algebra \(M\) and
every positive integer \(k\). Sinclair and Smith trace this
self-coefficient formulation to 1967~\cite[p.~3376]{SinclairSmith1998}.
Our main result settles the higher degrees.

\begin{theorem}\label{thm:main}
Let \(M\) be a complex von Neumann algebra and let \(k\geq2\).
Every \(f\in C_b^k(M,M)\) satisfying \(df=0\) has the form
\(f=dg\) for some \(g\in C_b^{k-1}(M,M)\).
Consequently \(H_b^k(M,M)=0\).
\end{theorem}

Together with the inner-derivation theorem,
Theorem~\ref{thm:main} resolves the Kadison--Ringrose cohomology
conjecture positively.

\subsection{Prior work}

Kadison and Ringrose proved vanishing for type~I and hyperfinite von
Neumann algebras~\cite[Theorem~4.4]{KadisonRingrose1971}
\cite[Theorem~3.1]{KadisonRingrose1971II}. Their extended-cobounding
theorem also gives a primitive for an $M$-valued cocycle on a general
concrete von Neumann algebra $M\subset B(H)$ when the primitive may
take values in
$B(H)$~\cite[Theorem~2.4]{KadisonRingrose1971II}; keeping the values in
the original algebra is the essential self-coefficient problem here.
Johnson, Kadison, and Ringrose then identified continuous cohomology
with separately normal cohomology~\cite[Theorem~5.6]{JKR1972}.
The completely bounded methods of Christensen, Effros, and
Sinclair~\cite{CES1987} established vanishing whenever the type
$\mathrm{II}_1$ central summand absorbs the hyperfinite
$\mathrm{II}_1$ factor, including all algebras without a type
$\mathrm{II}_1$ summand; see the precise restatement
in~\cite[p.~636, equation~(1.1)]{CPSS2003}.

The remaining finite-algebra problem has several substantial positive
cases. Sinclair and Smith proved vanishing for algebras with separable
predual and a Cartan subalgebra~\cite{SinclairSmith1998Cartan};
Cameron removed the separability hypothesis for
type~$\mathrm{II}_1$ factors with a Cartan maximal abelian
subalgebra~\cite[Theorem~3.2]{Cameron2009}.
Christensen, Pop, Sinclair, and Smith proved vanishing for all
type~$\mathrm{II}_1$ factors with property~$\Gamma$, including
nonseparable ones~\cite[Theorem~7.2]{CPSS2003}.
Qian and Shen extended this direction to type~$\mathrm{II}_1$
algebras with separable predual and their generalized
property~$\Gamma$~\cite[Theorem~6.4]{QianShen2015}.
In degree two, Pop and Smith proved vanishing for the tensor product
of any two type~$\mathrm{II}_1$ von Neumann algebras, with no
separability assumption~\cite[Theorem~4.5]{PopSmith2010}.
Theorem~\ref{thm:main} removes the structural restrictions on the
finite algebra in every degree $k\geq2$.

The Cartan argument of Sinclair and Smith also provides a useful
methodological comparison. They show that complete boundedness in the
last cocycle variable suffices. They use a projection of Christensen
and Sinclair~\cite{ChristensenSinclair1995} from completely bounded
linear maps onto right module maps, in the form recalled
in~\cite[Theorem~3.4]{SinclairSmith1998}, to obtain a primitive
\cite[Section~4]{SinclairSmith1998}.
Thus the passage from averaging one variable to a module identity
already has a central role in the earlier theory.

\subsection{The new ingredients}

We obtain the needed module identity by averaging over a single
unitary random walk. Averaging one cocycle input first gives only a
harmonic map; the main new input is a Liouville theorem that makes
such a map commute with right multiplication.
Suppose that \(M\) is a type \(\mathrm{II}_1\) algebra with separable
predual and faithful normal tracial state \(\tau\), and write
\(\norm{x}_2=\tau(x^*x)^{1/2}\). Let $\U(M)$ denote its unitary group.
We construct a symmetric, countably supported probability measure
\(\mu\) on \(\U(M)\), with a \(2\)-dense set of atoms, such that
\begin{equation}\label{eq:intro-harmonic}
T(x)=\E_{U\sim\mu}\bigl[T(xU)U^*\bigr]
\quad(x\in M)
\qquad\Longrightarrow\qquad
T(x)=T(1)x.
\end{equation}
Here \(T:M\to M\) is any bounded linear map that is \(2\)-norm
continuous on each operator-norm ball. No positivity or
complete-boundedness assumption is imposed.
The measure is fixed before the harmonic map is chosen.

For a separately normal cocycle, Akemann's continuity theorem
\cite[Corollary~III.10]{Akemann1967} supplies the stated continuity in
the last input. A uniform continuity modulus survives the averaging
limit, even though that limit need not remain normal, and also passes
to the tracial ultrapowers used in the Liouville proof.

The proof of \eqref{eq:intro-harmonic} combines probabilistic and
algebraic arguments. The finite free configurations come from Popa's
construction of free-independent sequences~\cite{Popa1995}, in the
formulation of his later ultraproduct independence
theorem~\cite{Popa2014}. These give finite unitary configurations with
prescribed small word traces. A hierarchy of increasingly rare,
increasingly large sets of such unitaries gives a single random walk
whose long independent products approach free Haar joint moments
along a subsequence. Only moment convergence is required.

Second, we use the classical harmonic-martingale construction
\cite[p.~415]{Furstenberg1963} to obtain an \(L^2(M,\tau)\)-valued
martingale. A nonconstant limit would produce free Haar unitaries
\(g_1,g_2,\ldots\) in a tracial ultrapower and distinct \(a,b\in M\).
The coordinatewise lift of \(T\) to that ultrapower sends each
nonempty signed word \(w\) with pairwise distinct unsigned indices
to \(aw\) or \(bw\), according to the sign of its first letter;
here \(a,b\) are identified with their constant images.
We prove that these identities are impossible for a bounded
map with the stated continuity.
The argument averages free generators, discards repeated surviving
indices using Haagerup's free-group length estimate~\cite{Haagerup1979},
and passes to a unitary
polar part. Positive and negative powers of that unitary would then
carry different left multipliers. Uniformly bounded sine sums and a
spectral partition force the multipliers to agree.
This rigidity statement, proved in
Theorem~\ref{thm:rigidity}, is independent of the cocycle problem.

Finally, in this tracial type \(\mathrm{II}_1\) setting, for a
separately normal cocycle \(f\) we take a pointwise ultraweak cluster
point of the Ces\`aro averages in its last input.
The Liouville theorem identifies the resulting fixed map as a
right module map in that input. Averaging the original cocycle equation
then gives a primitive \(g\) with
\(\norm{g}\leq\norm{f}\).
Together with compactness, this uniform estimate removes the
separability assumption. We then partition the type
\(\mathrm{II}_1\) summand into central pieces admitting faithful
normal tracial states. Lemma~\ref{lem:central-homotopy} corrects
mixed inputs from distinct pieces, and uniform bounds permit their
assembly. The classical vanishing theorem handles the complementary
summand.

\subsection{Conventions and classical reductions}

All von Neumann algebras and maps are complex, and subalgebras are
unital unless otherwise stated. A finite algebra is called
type \(\mathrm{II}_1\) here when it has no nonzero type~I central
summand; its center need not be trivial.
A map \(f:M^k\to M\) is \emph{separately normal} if it is
ultraweakly continuous in each input with the others fixed.
For a finite algebra with faithful normal tracial state, the
\(2\)-norm on bounded sets determines the strong-star topology.
The inherited traces on subalgebras and tracial ultrapowers will
usually also be denoted by \(\tau\).

We use the following established cohomology results.
\begin{samepage}
\begin{proposition}[Classical reductions]\label{prop:classical}
Let \(M\) be a von Neumann algebra and \(k\geq2\).
\begin{enumerate}[label=\textup{(\roman*)}]
\item If \(f\in C_b^k(M,M)\) is a cocycle, there are a separately
normal bounded cocycle \(f_0\) and a bounded \((k-1)\)-cochain \(b\)
such that \(f=f_0+db\).
\item If \(M\) has no type \(\mathrm{II}_1\) central summand, then
\(H_b^k(M,M)=0\).
\end{enumerate}
\end{proposition}
\end{samepage}

Part~\textup{(i)} is the Johnson--Kadison--Ringrose normal-reduction
theorem~\cite[Lemma~5.4 and Theorem~5.6]{JKR1972}; its formulation for self-coefficients is recalled
in \cite[Section~2]{CPSS2003}.
Part~\textup{(ii)} combines the classical type~I and properly infinite
vanishing results; see \cite[Theorem~4.4]{KadisonRingrose1971},
\cite{CES1987}, and \cite[p.~636, equation~(1.1)]{CPSS2003}.
Both statements concern ordinary continuous cohomology and the actual
image of \(d\). These are the only cohomological inputs in the proof.

Section~\ref{sec:walk} constructs the walk.
Section~\ref{sec:liouville} proves its Liouville property using
Theorem~\ref{thm:rigidity}, proved in Section~\ref{sec:rigidity}.
Section~\ref{sec:cohomology} constructs bounded primitives and
completes the reductions to arbitrary \(M\).
